\addcontentsline{toc}{section}{Resumen}
\begin{abstract}
Este trabajo da continuidad al estudio del Problema del Clique Máximo (MCP) iniciado en el Informe 1, donde se resolvió de forma exacta mediante un modelo de Programación Lineal Entera (ILP) resuelto por Ramificación y Acotamiento. En esta segunda etapa se aborda el mismo problema mediante una metaheurística de trayectoria: Búsqueda Tabú con reinicio por perturbación (estilo Búsqueda Local Iterada), construida sobre la representación de soluciones y el operador de vecindad Add/Drop/Swap especificados en la Tarea 1 de la asignatura. El algoritmo se implementa en Python y se evalúa mediante 30 corridas independientes sobre dos instancias del benchmark DIMACS: \textit{keller4} (171 vértices, $\omega=11$), la misma utilizada en el Informe 1, y \textit{C125.9} (125 vértices, densidad 0.90, $\omega=34$), una instancia más densa sobre la cual un Branch \& Bound exacto no logró probar optimalidad dentro de un límite de 300 segundos. Los resultados muestran que la metaheurística alcanza el óptimo conocido en el 93.3\,\% de las corridas en \textit{keller4} y en el 96.7\,\% en \textit{C125.9}, con un tiempo medio hasta la mejor solución de apenas 0.020--0.024 segundos, tres órdenes de magnitud menor que el tiempo del método exacto en ambos casos. El contraste directo entre ambas soluciones implementadas evidencia el compromiso fundamental entre garantía de optimalidad y escalabilidad: el método exacto certifica el óptimo pero su costo crece exponencialmente con la dificultad de la instancia, mientras que la metaheurística entrega soluciones de calidad prácticamente óptima en tiempos casi constantes, sin certificado de optimalidad. Se concluye que, para instancias donde el enfoque exacto se vuelve computacionalmente inviable, la Búsqueda Tabú con reinicio constituye una alternativa efectiva y de bajo costo computacional.

\vspace{0.4cm}
\noindent \textbf{Palabras clave:} Clique Máximo, Búsqueda Tabú, Búsqueda Local Iterada, Metaheurísticas, Optimización Combinatoria, Benchmark DIMACS, Contraste de métodos.
\end{abstract}
